\scanpage{34}
\deleted{entiers $i,j,\ell$, \unsure{donnés} par leurs développements $p$-adiques}

Soit $r\geq0$ tel que $p^r>i+j$, soit \deleted{$\unsure{\ell}$}
\[
\ell=p^r-1-(i+j)
\]
\deleted{Je dis que}
de sorte que l'on a
\[
i+j+\ell=p^r-1=\sum_{\alpha=0}^{r-1}(p-1)p^\alpha
\]
\inserted{$(i+j)+\ell=i+(j+\ell)$}
On sait alors, par le critère habituel, que l'on a $c_{i+j,\ell}\not\equiv0\ (p)$, et pour la $=$ raison on a $c_{i,j+\ell}\not\equiv0\ (p)$

cqfd

\uline{NB} Le lemme commun qui a servi dans les lemmes 1 et 2 est le suivant:
Si $\mu,\nu\geq0$ sont tels que
\[
\mu+\nu=p^r-1\qquad(r\geq0\text{ convenable})
\]
alors
\[
c_{\mu,\nu}\not\equiv0\ (p)
\]
\deleted{(plus précisément, $\unsure{c_{\mu,\nu}}$}
\margin{\unsure{5|S}}
\mspage{33}
